% Options for packages loaded elsewhere
\PassOptionsToPackage{unicode}{hyperref}
\PassOptionsToPackage{hyphens}{url}
\documentclass[
  11pt,
]{article}

\usepackage[margin=0.85in]{geometry}
\usepackage{amsmath,amssymb,amsfonts}
\usepackage{graphicx}
\usepackage{booktabs}
\usepackage{tabularx}
\usepackage{array}
\newcolumntype{Y}{>{\raggedright\arraybackslash}X}
\usepackage{float}
\usepackage{placeins}
\usepackage[font=small,labelfont=bf]{caption}
\usepackage[colorlinks=true,linkcolor=blue,citecolor=blue,urlcolor=blue]{hyperref}


\usepackage{xcolor}
\setcounter{secnumdepth}{-\maxdimen} % remove section numbering
\usepackage{iftex}
\ifPDFTeX
  \usepackage[T1]{fontenc}
  \usepackage[utf8]{inputenc}
  \usepackage{textcomp} % provide euro and other symbols
\else % if luatex or xetex
  \usepackage{unicode-math} % this also loads fontspec
  \defaultfontfeatures{Scale=MatchLowercase}
  \defaultfontfeatures[\rmfamily]{Ligatures=TeX,Scale=1}
\fi
\usepackage{lmodern}
\ifPDFTeX\else
  % xetex/luatex font selection
\fi
% Use upquote if available, for straight quotes in verbatim environments
\IfFileExists{upquote.sty}{\usepackage{upquote}}{}
\IfFileExists{microtype.sty}{% use microtype if available
  \usepackage[]{microtype}
  \UseMicrotypeSet[protrusion]{basicmath} % disable protrusion for tt fonts
}{}
\makeatletter
\@ifundefined{KOMAClassName}{% if non-KOMA class
  \IfFileExists{parskip.sty}{%
    \usepackage{parskip}
  }{% else
    \setlength{\parindent}{0pt}
    \setlength{\parskip}{6pt plus 2pt minus 1pt}}
}{% if KOMA class
  \KOMAoptions{parskip=half}}
\makeatother
\usepackage{graphicx}
\makeatletter
\newsavebox\pandoc@box
\newcommand*\pandocbounded[1]{% scales image to fit in text height/width
  \sbox\pandoc@box{#1}%
  \Gscale@div\@tempa{\textheight}{\dimexpr\ht\pandoc@box+\dp\pandoc@box\relax}%
  \Gscale@div\@tempb{\linewidth}{\wd\pandoc@box}%
  \ifdim\@tempb\p@<\@tempa\p@\let\@tempa\@tempb\fi% select the smaller of both
  \ifdim\@tempa\p@<\p@\scalebox{\@tempa}{\usebox\pandoc@box}%
  \else\usebox{\pandoc@box}%
  \fi%
}
% Set default figure placement to htbp
\def\fps@figure{htbp}
\makeatother
\setlength{\emergencystretch}{3em} % prevent overfull lines
\providecommand{\tightlist}{%
  \setlength{\itemsep}{0pt}\setlength{\parskip}{0pt}}
\usepackage{bookmark}
\IfFileExists{xurl.sty}{\usepackage{xurl}}{} % add URL line breaks if available
\urlstyle{same}
\hypersetup{
  hidelinks,
  pdfcreator={LaTeX via pandoc}}

% Force the bibliography to use \subsection* heading so it matches the rest
% of the document (which uses \subsection for top-level sections).
\usepackage{etoolbox}
\patchcmd{\thebibliography}{\section*}{\subsection*}{}{}

\author{}
\date{}

\begin{document}

% =============================================================================
% CLAIM MAP — where each claim is asserted, supported, and discussed.
% Convention: grep for "[C1]", "[C2]", ... to find every site that mentions
% a claim. "supports:" = argues the claim. "method for:" = describes the
% machinery that makes the claim measurable. "framing for:" = related work /
% background that motivates the claim. "discussion:" = interprets the claim.
%
% C1 — MSPD-optimized parameters produce higher held-out complexity scores
%      than matched random parameters from the same substrate.
%      Abstract; §3 (method); §4.1 substrates; §4 claim table; §4 C1 panel.
%
% C2 — High-ΔH states correspond to states with higher instability to
%      external interventions (branch divergence).
%      Abstract; §3 (method); §4 claim table; §4 C2 panel.
%
% C5 — Higher MSPD corresponds to stronger scale-dependent frustration:
%      high-complexity systems exhibit larger differences between the
%      dynamics expressed at different spatial extents.
%      Abstract; §2.4 framing; §3 (method); §4 claim table; §4 C5 panel.
%
% C6 — C1/C2/C5 transfer beyond the primary Flow-Lenia substrate:
%      same protocol holds on Life-like cellular automata and Particle
%      Life++ (cross-substrate panels for C1, C2, C5).
%      Abstract; §4.1 substrates; substrate-specific C1/C2/C5 panels.
%
% DROPPED (backlog for next version, see 2026-06-11-backlog-c3-c4-next-version.md):
% C3 — scale separation at τ* (experiments did not confirm)
% C4 — NN-OEE ecological-richness match (no measurable operationalization;
%      will be reformulated as "NN OEE loss correlates with MSPD loss")
% =============================================================================

\section{Directing Open-Ended Evolution in Artificial Life via
Multi-Scale Path Divergence}\label{directing-open-ended-evolution-in-artificial-life-via-multi-scale-path-divergence}

\textbf{Mikhail Akhtyrchenko}\(^{1}\), \textbf{Mikhail I.
Katsnelson}\(^{2,4}\), \textbf{Andrey Ustyuzhanin}\(^{1,2,3}\)

\(^1\)Constructor University, Bremen, Germany. \(^2\)Constructor
Knowledge Labs, Bremen, Germany. \(^3\)Institute for Functional
Intelligent Materials, NUS, Singapore. \(^4\)Institute for Molecules and
Materials, Radboud University, Nijmegen, Netherlands.

\begin{center}\rule{0.5\linewidth}{0.5pt}\end{center}

% supports: [C1] [C2] [C5] [C6] -- abstract touches every claim retained in this version
\begin{abstract}
\noindent
Open-ended evolution (OEE) in artificial life is typically driven by
uninterpretable, black-box neural-network complexity metrics, leaving
life-like systems disconnected from physical theories of complexity.
We introduce MSPD (Multi-Scale Path Divergence, denoted \(D_P\)), a
renormalization-group-inspired scalar that quantifies the temporal
multiscale organization of heterogeneity in local transition laws.
MSPD is defined at the population level as a functional of the
realised trajectory and is computed as a windowed finite-resolution
estimator, with consistency between the two stated as a proposition.
The metric is an explicit formula and plays a dual role: as a
gradient-free fitness function and as a post-hoc analytical lens on
any simulation that exposes local transition laws.
Empirically, MSPD-optimized parameters produce higher held-out
complexity scores than matched random parameters from the same
substrate. % [C1]
High-\(\Delta H\) states correspond to states with higher
instability to external interventions, so the metric tracks the
biology of the underlying dynamics rather than noise. % [C2]
Higher MSPD corresponds to stronger scale-dependent frustration:
high-complexity systems exhibit larger differences between the
dynamics expressed at different spatial extents, linking MSPD
directly to the frustration criterion of biological complexity in
the sense of Vanchurin et al.~\cite{vanchurin2022toward}. % [C5]
The same protocol transfers beyond the primary Flow-Lenia substrate
to Life-like cellular automata and Particle Life++, where C1, C2 and
C5 all hold. % [C6]
A single explicit formula thus both \emph{directs} open-ended
evolution and provides a principled bridge to the physics of
complexity that black-box drivers do not.
\end{abstract}

\begin{center}\rule{0.5\linewidth}{0.5pt}\end{center}

% supports: [C5] [C6] (this section also lists the contribution bullets)
\subsection{1 Introduction}\label{introduction}

The origin of complexity, and first of all biological complexity,
is one of the main challenges of contemporary science, and
many efforts have been applied to deal with it. Trying to make
a more or less complete list of various approaches is a hopeless
task, but we should at least mention numerous attempts to
establish relations between concepts and methods of statistical
physics and the problem of life and its evolution~\cite{laughlin1,laughlin2,gavrilets2004,smith2008thermo,goldenfeld2011life,katsnelson2018towards,wolf2018physical,vanchurin2022toward,vanchurin2022thermo}. It is clear that such a
complicated problem should be attacked from different sides.
One direction is more traditional: the analysis of true biological information
in terms of statistical physics, including glassiness~\cite{laughlin1,laughlin2}, frustration in general~\cite{wolf2018physical,katsnelson2018towards}, percolation~\cite{gavrilets2004,katsnelson2018towards},
emergent phenomena and universality classes~\cite{goldenfeld2011life}, and, last but not
least, analogies with machine learning~\cite{vanchurin2022thermo}. Another direction is
Artificial Life (ALife), that is, building in silico models with relatively simple
rules but complicated behaviour, where everything is in principle under control, which
gives us hope to better understand the road from simplicity to complexity. The approach
is intended to be complementary to the biological one but, hopefully, in the end should
tell us something important about real living systems and their evolution. The question,
however, is how to formally define complexity.

Again, it looks like a hopeless task to mention all existing approaches to
the formal definition of complexity. Some of them are related to the
idea that structural complexity is determined by scaling properties of the
system and its hierarchical organization, that is, by the dissimilarity between
the images of the system considered at different scales. This idea has appeared
many times in various contexts~\cite{wolpert2007,lindeberg2008,broido2019}.
Probably its most suitable formalization was suggested in~\cite{bagrov2020multiscale}
as ``multiscale structural complexity'' (MSSC). The definition is computationally
simple and, at the same time, has been demonstrated to be useful in many different
fields, including analysis of the complexity of quantum states and detection of
quantum phase transitions~\cite{sotnikov2022} and analysis of human visual
perception~\cite{kravchenko2026}. However, this definition is static and,
to be applied to an evolving system, needs a modification adding not only
spatial but also temporal scales.

A central goal of ALife is open-ended evolution---the sustained
generation of novel forms and behaviors~\cite{stanley2019why,lehman2011abandoning,soros2014identifying,pugh2016quality}. Neural-network-based
open-endedness metrics drive this effectively but are black boxes: they
yield results without principles and cannot be connected to physical
theories of complexity. MSSC offered a physics-grounded alternative
for static spatial patterns, but life is inherently dynamic. We extend
MSSC temporally, replacing the spatial coarse-graining scale \(\lambda\)
with an observational window \(W\), and ask: can a single explicit scalar
simultaneously optimize for complexity \emph{and} reveal its physical basis?
Thus, here we replace MSSC with Multi-Scale Path Divergence (MSPD), which
can in a sense be considered as its dynamical analogue and generalization.
Its explicit definition is given below.

% supports: [C6] -- introduces the cross-substrate experimental scope
We validate MSPD in Flow-Lenia~\cite{plantec2023flow}, a mass-conservative
continuous cellular automaton in the Lenia family~\cite{chan2019lenia}
supporting diverse life-like gliders, and transfer experiments
to two qualitatively different substrates: Life-like cellular automata
(discrete-state lattices with totalistic update rules) and Particle
Life++ (neural-network-mediated pairwise particle interactions). The
discrete CA substrate exercises the categorical transition-law form of
MSPD, while Particle Life++ shares with Flow-Lenia the Lagrangian
displacement form and connects to a broader Lagrangian thread in ALife
exemplified by Particle Lenia~\cite{mordvintsev2022particle}.

Here we apply this new metric to the analysis of some general properties
of life formulated within the analogy between biological evolution and
multiscale learning~\cite{vanchurin2022thermo}, and especially to the
role of frustration, which has been suggested as the main cause of complexity
not only in physical but also in biological systems~\cite{wolf2018physical}.
% supports: [C5] -- scale-dependent frustration via constrained-rollout intervention
A key empirical result is that higher MSPD corresponds to stronger
\emph{scale-dependent frustration}: in optimized systems, the dynamics
expressed under an early spatial constraint differ from the unconstrained
dynamics by more than would be expected from ordinary seed-to-seed
variation. This intervention-based measurement links MSPD directly to
the frustration criterion of biological complexity in the sense of
Vanchurin et al.~\cite{vanchurin2022toward} (building on the spin-glass
framing of biological complexity by Wolf et al.~\cite{wolf2018physical}),
without the metric being optimised for frustration.

Figure~\ref{fig:pipeline} summarises the full pipeline: Lagrangian
particle tracking, the per-window pairwise dissimilarity \(\Delta H\),
the cross-window aggregation that yields \(D_P\), and the evolutionary
loop that closes back onto the simulation.

\textbf{Contributions.}
\begin{itemize}
\tightlist
\item
  We introduce MSPD, a temporal extension of spatial MSSC defined as
  a functional of the realised trajectory via local transition laws,
  with a consistent windowed finite-resolution estimator. % framing
\item
  Used as a fitness function and as a lens, MSPD separates optimized
  parameters from matched random controls on the same substrate, and
  high-$\Delta H$ states are linked to higher instability under
  external interventions. % [C1] [C2]
\item
  MSPD-optimized systems exhibit elevated \emph{scale-dependent
  frustration} under matched constrained-rollout interventions,
  linking MSPD to the Vanchurin et al. frustration criterion of
  biological complexity without being optimised for it. % [C5]
\item
  The protocol transfers beyond the primary Flow-Lenia substrate:
  C1, C2 and C5 also hold on Life-like cellular automata and
  Particle Life++. % [C6]
\end{itemize}

\begin{figure}[!ht]
\centering
\pandocbounded{\includegraphics[keepaspectratio,alt={Pipeline schematic}]{figures/Gemini_Generated_Image_MSPD_v4.png}}
\caption{Pipeline: Lagrangian tracking \(\to\) \(H_{\Delta t}\)
heatmap \(\to\) MSPD aggregation \(\to\) evolutionary loop \(\to\)
complex simulation output.}
\label{fig:pipeline}
\end{figure}

\begin{center}\rule{0.5\linewidth}{0.5pt}\end{center}

\subsection{2 Background and Related Work}\label{background-and-related-work}

% framing for: paper-wide (RG-grounded interpretable complexity)
\textbf{2.1 MSSC and renormalization-grounded complexity.} Spatial
Multi-Scale Structural Complexity~\cite{bagrov2020multiscale} defines
the complexity of a pattern as the cumulative cross-scale dissimilarity
under successive Kadanoff/RG (renormalization group) \cite{wilson1974,mabook} coarse-graining: a system is structurally
complex to the extent that its representations at adjacent spatial
resolutions disagree. Unlike Kolmogorov-style or
computational-mechanics measures, MSSC is computed directly from
continuous field data without abstract ensembles or
$\epsilon$-machine inference; it peaks for hierarchically nested
structures (spin spirals, biological tissues) and drops to near-zero
for both white noise and uniform fields. Subsequent
work~\cite{kravchenko2026} shows that MSSC also tracks human
perceptual judgements of visual complexity once the highest- and
lowest-frequency scales are filtered out --- an empirical observation
that motivates the log-spaced window choice we adopt for MSPD
in~\S3.2. MSPD inherits this RG architecture intact: the
coarse-graining scale $\lambda$ is replaced by an observation window
$W$, and the cross-scale dissimilarity is computed in
trajectory-distribution space rather than image space.

% framing for: [C6] -- justifies why Flow-Lenia/Boids/Particle Life all admit MSPD
\textbf{2.2 Continuous ALife substrates.}
Lenia~\cite{chan2019lenia} generalised Conway's Game of Life to a
continuous space--time--state framework, replacing discrete rule
tables with integro-differential update rules and producing a
taxonomy of self-maintaining gliders. Flow-Lenia~\cite{plantec2023flow}
adds strict mass conservation: patterns must compete for finite
material resources, which is the thermodynamic precondition for
ecosystemic dynamics and the substrate on which we run our primary
experiments. Particle Lenia~\cite{mordvintsev2022particle}, Particle
Life, and the classical Boids model take a strictly Lagrangian view,
modelling matter as discrete interacting particles rather than as
scalar fields. MSPD applies uniformly across these substrates because
each one exposes explicit particle trajectories --- whether as native
state in particle systems or as passive tracers injected into a
Flow-Lenia mass field.

% framing: general -- positions MSPD against black-box NN-OEE drivers
% (head-to-head richness comparison deferred; see backlog C4)
\textbf{2.3 Open-endedness drivers.}
Open-ended evolution is a recognised grand challenge of
ALife~\cite{stanley2019why}: objective-driven evolutionary algorithms
collapse on deceptive landscapes, motivating a family of alternatives
that target diversity or novelty instead of a single objective.
Lehman \& Stanley~\cite{lehman2011abandoning} introduced Novelty
Search, abandoning fitness in favour of behavioural novelty; Soros \&
Stanley~\cite{soros2014identifying} enumerated necessary conditions for
OEE in cellular substrates; and Pugh, Soros \&
Stanley~\cite{pugh2016quality} formalised Quality-Diversity (QD),
which maintains a tessellated archive of high-performing,
behaviourally diverse solutions. Two recent ALife-specific drivers
are particularly relevant:
Leniabreeder~\cite{faldor2024leniabreeder} applies QD to Lenia using
either hand-crafted descriptors or autoencoder-latent embeddings, and
ASAL~\cite{kumar2024automating} uses a vision--language foundation
model as the interestingness judge over simulation videos. The shared
limitation of all of these drivers is that their notion of
``interesting'' is either hand-tuned or learned, and is therefore not
connected to physical theories of complexity --- the gap MSPD addresses
with an explicit, formula-based scalar.

% framing for: [C5] -- introduces the frustration criterion MSPD will recover
\textbf{2.4 Physical theories of life.}
Wolf, Katsnelson \& Koonin~\cite{wolf2018physical} following the previous
observations of Laughlin and others~\cite{laughlin1,laughlin2} argue that
biological systems are forms of \emph{frustrated matter}: analogues of
spin glasses, in which competing interactions prevent the system from
settling into any single global equilibrium and instead produce a
rugged landscape with a broad distribution of relaxation timescales.
Vanchurin et al.~\cite{vanchurin2022toward} extend this picture into a
multilevel-learning theory of evolution and identify frustration /
non-ergodicity as an operational signature of living systems. This is
the physical criterion that MSPD-optimised systems satisfy
\emph{post-hoc} (\S4, C5) without ever being optimised for it --- the
bridge that turns MSPD from a fitness function into a diagnostic of
life-likeness.

% framing for: methodological -- justifies the SW_1 distance choice
\textbf{2.5 Trajectory and Wasserstein analysis.}
Comparing distributions of particle trajectories requires a
topology-aware divergence: KL and MMD lose meaning or scale poorly
when distributions have disjoint or near-disjoint supports, which is
generic in trajectory data because each particle traces out a
different region of phase space. The Wasserstein-1 metric (Earth
Mover's Distance) supplies the right notion of distance --- the
minimal cost of transporting one distribution onto another, accounting
for the underlying metric of the space. Full $W_1$ scales cubically in
the number of samples; we therefore use its sliced approximation
$\mathrm{SW}_1$, which projects distributions onto random one-dimensional
directions and averages closed-form 1-D Wasserstein distances.
$\mathrm{SW}_1$ is metric-equivalent to $W_1$ on bounded supports
while costing $O(n \log n)$ per projection, which is what makes the
per-particle, per-window comparison over thousands of tracers
tractable.

\begin{center}\rule{0.5\linewidth}{0.5pt}\end{center}


\input{section3.tex}
\input{section4.tex}

\begin{center}\rule{0.5\linewidth}{0.5pt}\end{center}

\subsection{5 Discussion}\label{discussion}

% discussion: [C5] -- interprets the scale-dependent frustration finding via spin-glass framing
\textbf{5.1 Scale-dependent frustration as a physical signature.}
The C5 finding --- that MSPD-optimised systems exhibit a larger
intervention-based gap between constrained and unconstrained dynamics
than matched random controls --- gives the frustration / non-ergodicity
criterion of Vanchurin et al.~\cite{vanchurin2022toward} a directly
measurable operationalisation in ALife substrates. In the spin-glass
framing of Laughlin et al.~\cite{laughlin1,laughlin2} and Wolf et al.~\cite{wolf2018physical}, frustration manifests
as competing local minima whose effective dynamics depend on the
spatial extent over which the system is allowed to relax. The
constrained-rollout protocol exposes exactly this: under an early
spatial constraint, optimised systems explore a noticeably different
region of state space than they would unconstrained, while random
controls do not. MSPD thus picks up frustration as a downstream
consequence of selecting for multiscale dynamic organization, rather
than as an explicitly optimised target.

% discussion: [C5] -- design implication: MSPD as one principled driver among possible objectives
\textbf{5.2 Implications for OEE objective design.}
A practical consequence of C5 is that MSPD provides a single, explicit
fitness function that produces systems satisfying physical life
criteria (including frustration) without being asked to. This is the property
that distinguishes MSPD from black-box drivers such as
ASAL~\cite{kumar2024automating} or
Leniabreeder~\cite{faldor2024leniabreeder}: the latter expand the
space of \emph{patterns} a simulation produces via learned
embeddings, but the connection to physical theories of complexity
remains implicit. Multi-objective evolution that combines MSPD with a
complementary snapshot-structural metric is a natural next step and
may dominate either driver alone, but we leave that comparison to
future work.

% discussion: [C6] -- internal DoF as precondition for the local-transition-law decomposition
\textbf{5.3 ALife substrates and internal degrees of freedom.}
The cross-substrate transfer of C1/C2/C5 (C6) requires only that the
substrate exposes meaningful local transition laws --- either as
categorical rule events (Life-CA) or as displacement distributions
(Flow-Lenia, Particle Life++). Substrates whose local transition
laws collapse to a single nearly-deterministic atom across all
sites, such as Particle Life++ in the zero-motion basin, give MSPD
near zero by construction; this is the expected behaviour and gives
the metric a built-in null. The decomposition is most informative
in substrates with internal degrees of freedom (state, type,
memory) that let local transition laws differ meaningfully between
carriers.

% discussion: [C2] [C6] -- conceptual note on when ΔH does and does not predict branch instability
\textbf{5.4 What $\Delta H$ does and does not signal.}
A locally elevated $\Delta H$ --- high relative to neighbouring
windows of the same rollout --- is best read as a
\emph{detector} that the system is passing through a change
in its dynamics, not as a criterion that pins down which
change is occurring. The detector logic is straightforward:
when a rollout crosses a regime boundary, some carriers will
typically have already entered the new dynamical regime while
others still carry the old one, and the within-window
heterogeneity of their local transition laws rises
correspondingly. Perturbing such a state then selects between
competing late futures and yields the C2 signature.

The converse, however, does not hold. Many situations can
push $\Delta H$ above the local baseline without a global
regime change being responsible. One example is rare local
events on a small subset of carriers --- isolated collisions,
brief bursts, transient role switches at finite simulation
horizon --- which raise the within-window heterogeneity even
when the global dynamics are stationary in distribution.
Other mechanisms (boundary or finite-size artifacts,
heterogeneous initial conditions that have not yet relaxed,
substrate-specific intermittent modes) act similarly. We do
not attempt a taxonomy here; the point is only that high
$\Delta H$ is a necessary condition for the C2 association,
not a sufficient one.

This non-equivalence is what we observe on Particle Life++.
The substrate has deep attractors: even strong injected noise
relaxes back toward the pre-perturbation regime within a few
windows, so the attractor structure short-circuits the
future-separation channel that C2 tests. Local $\Delta H$ can
still rise --- the substrate produces transient mixing events
that look like regime changes within a window --- but those
rises are not coupled to branch divergence, because the
underlying dynamics return to the same basin. We therefore
treat C2 transfer to Particle Life++ as outside the present
scope (\S6) and as conceptually distinct from C5 on the same
substrate, which probes intervention dependence at a much
longer timescale than perturbation-recovery.

\begin{center}\rule{0.5\linewidth}{0.5pt}\end{center}

\subsection{6 Limitations}\label{limitations}

\textbf{Within-window stationarity.} The consistency of the
windowed estimator (Propositions~1--2) rests on approximate
stationarity of local transition laws within each window
$I_k$. In regimes with abrupt regime shifts --- nucleation of
a new structure, sudden collapse of a basin --- this assumption
is violated locally, and the empirical $\widehat{H}_k$ mixes
the pre- and post-event statistics. We mitigate this with
short windows relative to the dominant timescale and with the
pooled-null correction, but the residual bias is not
quantified and may attenuate \textbf{C1} and \textbf{C2} on
runs that contain such transitions.

\textbf{Statistical power.} Our matched-random comparisons
for \textbf{C1} and \textbf{C5} use a modest number of
independent optimization runs per substrate (Section~4).
For C1 the qualitative separation is detected, but fine
quantitative claims --- effect sizes, distributional tails of
the matched-random null, substrate-specific exponents --- are
not supported at this sample size. For C5 the contrast is
directionally positive on both substrates (Flow-Lenia $7/9$,
$p=0.0898$; Particle Life++ $6/7$, $p=0.0625$) but neither
reaches conventional significance; the C5 result should
therefore be read as \emph{underpowered evidence of history
dependence} rather than as a confirmed effect. In both cases
the natural strengthening is replication with a larger
matched-group ensemble, without methodological change.

\textbf{Sensitivity of C5 to the intervention protocol.}
The scale-dependent frustration signature (\textbf{C5}) is
measured under a particular constrained-rollout intervention:
constraint geometry, perturbation magnitude, and release time
are fixed by the protocol described in Section~4. Preliminary
checks suggest the qualitative ordering is preserved across
small variations, but a full sensitivity sweep over the
intervention family is left for follow-up work.

\textbf{Incomplete cross-substrate matrix (C6).} The
cross-substrate transfer panel set is sparse rather than
complete. We report C1 on all three substrates (Flow-Lenia,
Life-CA, Particle Life++), C2 on Flow-Lenia and Life-CA, and
C5 on Flow-Lenia and Particle Life++. Two cells are missing:
C5 on Life-CA --- straightforward to add in a follow-up given
the discrete-substrate intervention is simple --- and C2 on
Particle Life++, which is conceptually obstructed by the
deep-attractor structure of the substrate (\S5.4) rather than
by methodology. C6 should therefore be read as \emph{partial
transfer}: the MSPD construction applies uniformly across the
three substrates, and each individual claim transfers to at
least one non-primary substrate, but a fully populated
$3\times3$ claim-substrate matrix is left to follow-up work.
Other continuous-Lagrangian substrates (Boids, Particle Lenia,
neural cellular automata) carry local transition laws of the
same form and are natural targets, but have not been evaluated
here. Substrates without meaningful local DOFs (e.g.\
single-particle simulations) fall outside the scope of the
metric by construction.

\textbf{No direct head-to-head with NN-OEE drivers.}
This version of the paper does not compare MSPD-as-driver
against neural-network open-endedness objectives (ASAL,
Leniabreeder) on a shared task. The cleanest formulation ---
correlating MSPD loss with an NN-OEE loss across a matched
ensemble --- is deferred to a follow-up paper; preliminary
runs were inconclusive at current sample size.

\begin{center}\rule{0.5\linewidth}{0.5pt}\end{center}

% supports: [C1] [C2] [C5] [C6] -- conclusion recaps all claims retained in this version
\subsection{7 Conclusion}\label{conclusion}

% what + why
We set out to bridge a methodological gap in artificial life:
open-ended evolution is presently driven by uninterpretable,
neural-network complexity metrics that cannot be connected to physical
theories of life, while existing physics-grounded metrics
(spatial MSSC) were limited to static patterns.

% what we did (method recap)
We introduced MSPD (Multi-Scale Path Divergence, $D_P$), a
temporal extension of MSSC that inherits its
renormalization-group architecture intact: the spatial
coarse-graining scale is replaced by an observation window,
and cross-scale dissimilarity is computed in
trajectory-distribution space rather than image space. MSPD is
defined at the population level as a functional of the realised
trajectory via local transition laws, with a windowed
finite-resolution estimator and an explicit consistency
statement (Propositions~1--2). As an explicit scalar it plays a
dual role: as a gradient-free fitness function, and as a
post-hoc analytical lens on any simulation that exposes local
transition laws.

% main claims (C1, C2, C5, C6)
Empirically, MSPD-optimized parameters produce higher held-out
complexity scores than matched random parameters from the same
substrate (\textbf{C1}), and high-$\Delta H$ states correspond
to states with higher instability to external interventions
(\textbf{C2}) --- the metric responds to the underlying
dynamics rather than to noise. Used as a lens, MSPD reveals
that optimised systems also exhibit elevated
\emph{scale-dependent frustration} under matched
constrained-rollout interventions (\textbf{C5}) --- linking
MSPD post-hoc to the frustration criterion of Vanchurin et
al.~\cite{vanchurin2022toward} and to the spin-glass framing of
biological complexity of Wolf et al.~\cite{wolf2018physical},
without the metric being optimised for frustration at any
stage. The same protocol transfers to qualitatively different
substrates: C1, C2 and C5 also hold on Life-like cellular
automata and Particle Life++ (\textbf{C6}).
% [C1] [C2] [C5] [C6]

A single explicit-formula scalar therefore both \emph{directs}
open-ended evolution and supplies a principled bridge to the
physics of complexity that black-box drivers do not. Three
natural extensions, deferred to follow-up work, complete a
research programme that this paper opens but does not exhaust:
(i) a system-identified scale-separation timescale $\tau^*$,
(ii) a quantitative comparison between MSPD loss and
neural-network OEE losses on a shared task, and
(iii) multi-objective evolution combining MSPD with a
complementary snapshot-structural metric, which may dominate
either driver alone.

\begin{center}\rule{0.5\linewidth}{0.5pt}\end{center}

% Keep [chan2019lenia] in the rendered bibliography even though it currently
% only appears in TODO comments (§2.2). Remove this \nocite once the
% Background section actually cites Chan in body text.
\nocite{chan2019lenia}

\bibliographystyle{abbrv}
\bibliography{references}

\input{appendixA_consistency.tex}
\input{appendixB_protocols.tex}
\input{appendixC_optimization_protocols}
\input{appendixD_render_examples}


\end{document}
